\subsection{Открытые табличные форматы}

Открытый табличный формат определяет, как из набора файлов данных, лежащих в
объектном хранилище, собирается логическая таблица. Минимальный набор обязанностей
табличного формата включает: ведение списка файлов, составляющих текущее состояние
таблицы; атомарную фиксацию изменений (добавление и удаление файлов одной
транзакцией); изоляцию снимков (читатель видит согласованное состояние на момент
начала чтения, не затрагиваемое параллельной записью); хранение схемы и её
эволюцию; статистики по файлам для пропуска нерелевантных данных; разбиение на
разделы. Рассмотрим основные реализации.

\textbf{Apache Iceberg}~\cite{iceberg-spec} ведёт метаданные в виде дерева: файл метаданных таблицы
ссылается на список манифестов (manifest list), каждый манифест перечисляет файлы
данных с их статистиками. Фиксация транзакции --- это атомарная замена указателя на
новый файл метаданных. Iceberg поддерживает скрытое разбиение на разделы (partition
evolution), эволюцию схемы без перезаписи данных, ветвление и теги снимков. Удаления
реализуются через позиционные и равенственные delete-файлы. Iceberg --- наиболее
широко принятый в индустрии формат, поддержанный всеми крупными движками.

\textbf{Delta Lake}~\cite{armbrust2020delta} использует журнал транзакций --- упорядоченную последовательность
JSON-файлов, каждый из которых описывает добавленные и удалённые файлы данных;
периодически журнал уплотняется в контрольную точку в формате Parquet. Delta
поддерживает обновления и удаления через перезапись затронутых файлов
(copy-on-write) либо через deletion vectors. Тесно интегрирован с Apache Spark.

\textbf{Apache Hudi}~\cite{hudi-docs} ориентирован на инкрементальную обработку и upsert по ключу
записи. Поддерживает два типа таблиц: copy-on-write (обновления порождают перезапись
файлов) и merge-on-read (обновления пишутся в журнальные файлы, состояние
вычисляется при чтении). Ведёт временную шкалу (timeline) операций и индекс,
сопоставляющий ключи записей файлам.

\textbf{Apache Paimon}~\cite{paimon-docs} (изначально Flink Table Store) выделяется тем, что состояние
таблицы с первичным ключом организовано как LSM-дерево, размещённое в объектном
хранилище. Свежие изменения сбрасываются в файлы уровня~0, фоновое слияние
перемещает их на нижележащие уровни; при чтении актуальная версия строки получается
объединением версий из всех уровней по полю последовательности. Метаданные ведутся
аналогично Iceberg: снимок ссылается на манифесты, манифесты --- на файлы данных со
статистиками. Такое устройство делает Paimon эффективным для потоковой записи,
захвата изменений и near-real-time-витрин --- сценариев, в которых таблица
постоянно обновляется небольшими порциями, а не перезаписывается целиком. Apache
Paimon поддерживает несколько файловых форматов хранения данных, выбираемых
параметром таблицы: Apache Parquet (по умолчанию), Apache ORC и Apache Avro;
добавление поддержки формата Vortex составляет практическую часть настоящей работы.

Для задач настоящей работы важно, что табличный формат и файловый формат --- разные
уровни абстракции: первый управляет составом и версиями набора файлов, второй ---
двоичным представлением данных внутри одного файла. Производительность чтения и
записи определяется обоими уровнями; в работе фиксируется табличный формат (Apache
Paimon) и варьируется файловый формат.
